% Source-derived chapter 08: Analytic $1$-parameter subgroups
% Original source: chabauty-coleman-bound-surfaces
\chapter{Analytic $1$-parameter subgroups}
\label{chabauty-coleman-bound-surfaces-chapter-08}
\source{chabauty-coleman-bound-surfaces}

\begin{remark}[Source section]
\label{chabauty-coleman-bound-surfaces-chapter-08-source}
\source{chabauty-coleman-bound-surfaces}
\source{chabauty-coleman-bound-surfaces}
This chapter reproduces the corresponding section of the retrieved
LaTeX source, including its definitions, statements, examples, and
proofs. It has not yet been translated into Lean.
\notready
\end{remark}

% The source section title was: Analytic $1$-parameter subgroups
\label{Sec1PS}
\source{chabauty-coleman-bound-surfaces}

Using the results on the logarithm and exponential maps from the previous section, in this section introduce analytic $1$-parameter subgroups of abelian varieties in the non-archimedian setting and study some fundamental properties. The naive idea is that a $1$-parameter subgroup of an abelian variety should be an analytic curve at the identity element which is locally a subgroup. However, it is more convenient for our purposes to take a different approach.

Again, we let $\Acal$ be an abelian variety over $\Spec R$ of relative dimension $n\ge 1$. We keep the notation introduced in Sections \ref{NotationLoc} and  \ref{SecLogExp}. We assume 
$$
p>\max\{\efrak +1 , \exp(\efrak/\exp(1))\}
$$ 
throughout this section, so that Lemma \ref{LemmaKeyLogExp} applies.

%%%%
%%%%
%%%%
\subsection{Definitions} If $G$ is an abelian $K$-analytic Lie group, the $K$-vector space of invariant differentials will be denoted by $\Omega(G)$. A morphism of abelian Lie groups $f:G_1\to G_2$ functorially induces a linear map on invariant differentials, which we denote by $f^\bullet:\Omega(G_2)\to \Omega(G_1)$.


A \emph{$1$-parameter subgroup} of $\Acal$ is a pair $\gamma=(\ttt,\uu)$ where $\ttt=(t_1,...,t_n)$ is a choice of $R$-local parameters at $e\in A(K)$ and $\uu\in K^n$ satisfies $|\uu|=1$.  


 The space of invariant differentials on the Lie group $\Tcal$ is $\Omega(\Tcal) = \bigoplus_{j=1}^n K\cdot dx_j$ and the space of invariant differentials on $A(K)$ is $ \Omega(A(K))=H^0(A,\Omega^1_{A/K})$. 
 
 Given $\ttt$  a choice of $R$-local parameters at $e\in A(K)$, the isomorphism of Lie groups $\widetilde{\Exp}^\ttt:\Tcal\to U_e$ (cf. Lemma \ref{LemmaKeyLogExp}) induces an isomorphism of $K$-vector spaces on invariant differentials 
\begin{equation}\label{Eqnzxc}
\source{chabauty-coleman-bound-surfaces}
(\widetilde{\Exp}^\ttt)^\bullet : H^0(A,\Omega^1_{A/K})\to \Omega(\Tcal).
\end{equation}
Given $\uu=(u_1,...,u_n)\in K^n$ with $|\uu|=1$, we define the $K$-linear map 
$$
\epsilon^\uu: \Omega(\Tcal)\to K
$$ 
by the rule $dx_j\mapsto u_j$. For a $1$-parameter subgroup $\gamma=(\ttt,\uu)$ we get a non-trivial $K$-linear map 
$$
\epsilon^\uu\circ (\widetilde{\Exp}^\ttt)^\bullet: H^0(A,\Omega^1_{A/K})\to K.
$$
Let us define $\Hcal(\gamma)=\ker(\epsilon^\uu\circ (\widetilde{\Exp}^\ttt)^\bullet)$ and note that it is a hyperplane in $H^0(A,\Omega^1_{A/K})$.

A $1$-parameter subgroup $\gamma=(\ttt,\uu)$ determines a morphism of Lie groups 
%%
\begin{equation}\label{Eqntildegamma}
\source{chabauty-coleman-bound-surfaces}
\tilde{\gamma}:\pfrak \to U_e, \quad \tilde{\gamma}(z)=\widetilde{\Exp}^\ttt(u_1z,...,u_nz)
\end{equation}
%%
where $z$ denotes the variable on the additive Lie group $\pfrak$. The space of invariant differentials on $\pfrak$ is $\Omega(\pfrak) = K\cdot dz$. One immediately verifies
%%
%%
\begin{lemma}
\source{chabauty-coleman-bound-surfaces}
\notready\label{LemmaHisKernel} The map $\tilde{\gamma}^\bullet : H^0(A,\Omega^1_{A/K})\to \Omega(\pfrak)$ induced by $\tilde{\gamma}$  is given by $\omega\mapsto (\epsilon^\uu\circ (\widetilde{\Exp}^\ttt)^\bullet)(\omega)\cdot dz$. In particular, $\ker(\tilde{\gamma}^\bullet)=\Hcal(\gamma)$.
\source{chabauty-coleman-bound-surfaces}
\end{lemma}
%%
%%
The \emph{image} of a $1$-parameter subgroup $\gamma$ is defined as $\im(\gamma)=\tilde{\gamma}(\pfrak)$. We observe that
%%
%%
\begin{lemma}
\source{chabauty-coleman-bound-surfaces}
\notready\label{LemmaSaturated} If $\gamma=(\ttt,\uu)$ is a $1$-parameter subgroup of $\Acal$, then $\im(\gamma)$ is a Lie subgroup of $U_e$. The image of $\im(\gamma)$ under the isomorphism $\widetilde{\Log}^\ttt:U_e\to \Tcal$ is $\pfrak\cdot \uu$, which is a closed, saturated $R$-submodule of $\Tcal$, free of rank $1$ over $R$, and generated by $\varpi\cdot \uu$.
\source{chabauty-coleman-bound-surfaces}
\end{lemma}
%%
%%


%%%%
%%%%
%%%%
\subsection{Equivalence} Given $\gamma$ and $\gamma'$ two $1$-parameter subgroups  of $\Acal$, we say that they are \emph{equivalent} if $\Hcal(\gamma)=\Hcal(\gamma')$. This is denoted by $\gamma\sim\gamma'$.

%%
%%
\begin{lemma}
\source{chabauty-coleman-bound-surfaces}
\notready \label{LemmaVaryu1PS} Let $\ttt$ be a choice of $R$-local parameters at $e$ and let $\uu,\uu'\in K^n$ with $|\uu|=|\uu'|=1$. Let $\gamma=(\ttt,\uu)$ and $\gamma'=(\ttt,\uu')$. We have $\gamma\sim \gamma'$ if and only if there is $\eta\in R^\times$ with $\uu'=\eta\cdot \uu$.
\source{chabauty-coleman-bound-surfaces}
\end{lemma}
%%
%%
\begin{proof} As \eqref{Eqnzxc} is an isomorphism, we have $\Hcal(\gamma)=\Hcal(\gamma')$ if and only if $\ker(\epsilon^\uu)=\ker(\epsilon^{\uu'})$. The latter holds if and only if $\uu'=\eta\cdot \uu$ for some $\eta\in K^\times$, in which case  $\eta\in R^\times$ as $|\uu|=|\uu'|=1$.
\end{proof}
%%
%%



%%
%%
\begin{lemma}
\source{chabauty-coleman-bound-surfaces}
\notready \label{LemmaEquiv1PS} Let $\gamma$ and $\gamma'$ be $1$-parameter subgroups of $\Acal$. The following are equivalent:
\source{chabauty-coleman-bound-surfaces}
\begin{itemize}
\item[(i)] $\gamma\sim\gamma'$.
\item[(ii)] $\im(\gamma)=\im(\gamma')$.
\item[(iii)] $\im(\gamma)$ and $\im(\gamma')$ have non-trivial intersection, i.e. $\{e\}\subsetneq \im(\gamma)\cap\im(\gamma')$.
\end{itemize}
\end{lemma}
%%
%%
\begin{proof}
Write $\gamma=(\ttt,\uu)$ and $\gamma'=(\ttt',\uu')$. Define $\delta:\Tcal\to \Tcal$ as $\delta=\widetilde{\Log}^\ttt\circ \widetilde{\Exp}^{\ttt'}$. Then $\delta$ is a Lie group automorphism of the additive Lie group $\Tcal$. Thus, $\delta$ is given by $\delta\in \GL_n(R)$ acting on $\Tcal=\pfrak\times...\times \pfrak$ on the left.

After these remarks, Lemma \ref{LemmaSaturated} shows that (ii) and (iii) are equivalent, and in fact, both are equivalent to the following condition: (iv) there exists $\eta\in R^\times$ with $\delta(\uu') = \eta\cdot \uu$ (note that as an element of $\GL_n(R)$, $\delta$ acts on $K^n$).

Let us show that (iv) is equivalent to (i). Note that the explicit definition of $\epsilon^{\uu}$ and the fact that $|\uu|=|\uu'|=1$ give that (iv) is equivalent to the condition $\ker(\epsilon^\uu) = \ker(\epsilon^{\delta(\uu')})$. Directly evaluating on $dx_i$ we see that $\epsilon^{\delta(\uu')}=\epsilon^{\uu'}\circ \delta^\bullet$, so, (iv) is equivalent to $\ker(\epsilon^\uu) = \ker(\epsilon^{\uu'}\circ \delta^\bullet)$. Using the fact that $\delta^\bullet\circ (\widetilde{\Exp}^{\ttt})^\bullet=(\widetilde{\Exp}^{\ttt}\circ \delta)^\bullet =( \widetilde{\Exp}^{\ttt'})^\bullet$ and that $( \widetilde{\Exp}^{\ttt})^\bullet$ is an isomorphism, we finally see that (iv) is equivalent to
$\ker(\epsilon^\uu\circ( \widetilde{\Exp}^{\ttt})^\bullet)=\ker(\epsilon^{\uu'}\circ( \widetilde{\Exp}^{\ttt'})^\bullet)$, which is exactly (i).
\end{proof}
%%
%%

\begin{lemma}
\source{chabauty-coleman-bound-surfaces}
\notready\label{LemmaH1PS} The rule $\gamma\mapsto \Hcal(\gamma)$ defines a bijection between equivalence classes of $1$-parameter subgroups of $\Acal$ and $K$-linear hyperplanes of $H^0(A,\Omega^1_{A/K})$.
\source{chabauty-coleman-bound-surfaces}
\end{lemma}
\begin{proof} By Lemma \ref{LemmaEquiv1PS}, it only remains to show that given any hyperplane $H$ of $H^0(A,\Omega^1_{A/K})$ there is a $1$-parameter subgroup $\gamma$ with $H=\Hcal(\gamma)$. Choose $R$-local parameters $\ttt$ at $e$. Let $H'=(\widetilde{\Exp}^\ttt)^\bullet(H)$ and note that it is a hyperplane in $\Omega(\Tcal)$. We can choose $\uu\in K^n$ with $|\uu|=1$ such that $\ker(\epsilon^\uu)=H'$, hence $H=\ker(\epsilon^\uu\circ (\widetilde{\Exp}^\ttt)^\bullet)$ as desired.
\end{proof}

%%
%%
\begin{lemma}
\source{chabauty-coleman-bound-surfaces}
\notready\label{Lemma1PSgen} Let $\xi\in U_e$ with $\xi\ne e$ and let $\ttt$ be a choice of $R$-local parameters at $e$. There exists $\uu\in K^n$ with $|\uu|=1$ such that $\xi \in \im(\gamma)$ for $\gamma=(\ttt,\uu)$. Furthermore, $\uu$ is unique for this $\xi$ and $\ttt$, up to multiplication by elements of $R^\times$.
\source{chabauty-coleman-bound-surfaces}
\end{lemma}
\begin{proof} Let $\vv=\widetilde{\Log}(\xi)\in \Tcal$. We have $\vv\ne 0$ because $\xi\ne e$ and $\widetilde{\Log}:U_e\to \Tcal$ is an isomorphism of Lie groups. Take $m\in \Z$ such that $|\varpi^m|=|\vv|$. We can choose $\uu=\varpi^{-m}\cdot \vv$. 

If $\gamma=(\ttt,\uu)$ and $\gamma'=(\ttt,\uu')$ satisfy that $\xi\in \im(\gamma)\cap\im(\gamma')$ then $\gamma\sim\gamma$ by Lemma \ref{LemmaEquiv1PS}, and we conclude by Lemma \ref{LemmaVaryu1PS}.
\end{proof}
%%%%
%%%%
%%%%
\subsection{Infinitesimal $1$-parameter subgroups} For an integer $m\ge 0$ we define 
\begin{itemize}
\item $\Ecal_m= R[z]/(z^{m+1})$, $\Vcal_m=\Spec \Ecal_m$, 
\item $E_m=\Ecal_m\otimes_R K$, $V_m=\Spec E_m$, 
\item $E'_m=\Ecal_m\otimes_R \kk$, $V'_m=\Spec E'_m$.
\end{itemize}

Let $\gamma=(\ttt,\uu)$ be a $1$-parameter subgroup of $\Acal$. The power series expansion of the Lie group morphism $\tilde{\gamma}:\pfrak\to U_e$ (see \eqref{Eqntildegamma}) with respect to $z$ and $\ttt$ induces a morphism of formal schemes $\widehat{\gamma}: \Spf K[[z]]\to  \Spf \widehat{\Ocal}_{A,e}$ determined by the following map on completed local rings:
%%
\begin{equation}\label{Eqnhatgamma}
\source{chabauty-coleman-bound-surfaces}
\widehat{\gamma}^\#: \widehat{\Ocal}_{A,e}\to K[[z]], \quad t_j\mapsto \Exp^\ttt_j(u_1z,...,u_nz).
\end{equation}
%%
For each $m\ge 0$, the map $\widehat{\gamma}$ induces a map $\widehat{\gamma}_m:V_m\to \Spf  \widehat{\Ocal}_{A,e}$ giving the morphism of $K$-schemes
$$
\tilde{\gamma}_m:V_m\to A
$$
supported at $e$. Explicitly, $\widehat{\gamma}_m$ induces 
%%
\begin{equation}\label{Eqninf1ps}
\source{chabauty-coleman-bound-surfaces}
\widehat{\gamma}_m^\#: \widehat{\Ocal}_{A,e}\to E_m, \quad t_j\mapsto \Exp^\ttt_j(u_1z,...,u_nz)\bmod z^{m+1}
\end{equation}
%%
on completed local rings, and the restriction to $\Ocal_{A,e}$ is the map $\tilde{\gamma}_m^\#:\Ocal_{A,e}\to E_m$ induced by $\tilde{\gamma}_m$. In summary, the following diagram commutes:
%%
\begin{equation}\label{Eqngammas}
\source{chabauty-coleman-bound-surfaces}
\begin{tikzcd}
\widehat{\Ocal}_{A,e}   \arrow{r}{\widehat{\gamma}^\#} \arrow{rd}{\widehat{\gamma}^\#_m}   &  K[[z]]  \arrow{d}  \\
\Ocal_{A,e}  \arrow[hook]{u}  \arrow[swap]{r}{\tilde{\gamma}^\#_m }&   E_m.
\end{tikzcd}
\end{equation}
%%
The morphism $\tilde{\gamma}_m$ is the $m$-th jet of $\tilde{\gamma}$, which might be thought of as an infinitesimal $1$-parameter subgroup. We remark that, although $\tilde{\gamma}$ is an analytic maps of Lie groups, the maps $\tilde{\gamma}_m:V_m\to A$ are scheme morphisms.

For each integer $h\ge 0$ it will be convenient to define the polynomial $P^{\ttt}_{j,h}(\xx)\in K[\xx]$ as the homogeneous part of degree $h$ of $\Exp^\ttt_j(\xx)$, which in the notation of Section \ref{SecLogExp2} can be written as
$$
P^{\ttt}_{j,h}(\xx)=\frac{1}{m!}\cdot \psi^{[h]}_j(\xx)=\sum_{\|\alpha\|=h} c_{j,\alpha}\cdot \xx^\alpha.
$$
In particular, $P^\ttt_{j,0}=0$. With this notation, the map $\widehat{\gamma}^\#$ from \eqref{Eqninf1ps} can be expressed as
%%
\begin{equation}\label{Eqninf1psP}
\source{chabauty-coleman-bound-surfaces}
\widehat{\gamma}^\#_m(t_j)=\sum_{h=0}^{m} P^\ttt_{j,h}(\uu)\cdot z^h\bmod z^{m+1}. 
\end{equation}
%%

%%
%%
\begin{lemma}[Integrality and reduction modulo $\pfrak$]
\source{chabauty-coleman-bound-surfaces}
\notready\label{LemmaTruncate1ps} Let $\gamma=(\ttt,\uu)$ be a $1$-parameter subgroup of $\Acal$. If $m<p$, the morphism $\tilde{\gamma}_m:V_m\to A$ extends to an $R$-morphism $\Vcal_m\to \Acal$ supported along the identity section $\sigma_e:\Spec R\to \Acal$. Hence, upon base change to $\kk$, it determines a $\kk$-morphism $V'_m\to A'$ supported at $e$. 
\source{chabauty-coleman-bound-surfaces}
\end{lemma}
\begin{proof} Lemma \ref{LemmaCvExp} shows that $|P^\ttt_{j,h}(\uu)|\le 1$ for each $h\le m$, since $m<p$. We conclude by applying this estimate to \eqref{Eqninf1psP} and the fact that $\ttt$ is a choice of $R$-local parameters for $\Acal$ along $e$.
\end{proof}
%%
%%
Given a $1$-parameter subgroup $\gamma$ and an integer $0\le m<p$, the morphisms provided by Lemma \ref{LemmaTruncate1ps} will be denoted by
$$
\tilde{\gamma}^R_m:\Vcal_m\to \Acal\quad \mbox{ and }\quad \tilde{\gamma}'_m : V'_m\to A'.
$$
%%
%%
\begin{lemma}[Closed immersions]
\source{chabauty-coleman-bound-surfaces}
\notready\label{LemmaClosedImm} Let $\gamma=(\ttt,\uu)$ be a $1$-parameter subgroup of $\Acal$. For each $m\ge 0$, the morphism $\tilde{\gamma}_m:V_m\to A$ is a closed immersion. If $m<p$, then the maps $\tilde{\gamma}^R_m:\Vcal_m\to \Acal$ and $\tilde{\gamma}'_m:V'_m\to A'$ are also closed immersions.
\source{chabauty-coleman-bound-surfaces}
\end{lemma}
\begin{proof} For $m=0$ we have $E_0=K$, $\Ecal_0=R$, and $E'_0=\kk$, so, the result is clear. Assume $m\ge 1$.

By Lemma \ref{LemmaLinearPart} we have $P^\ttt_{j,1}(\xx)=x_j$. Thus, $P^\ttt_{j,1}(\uu)=u_j$ and since $|\uu|=1$ we conclude that for some $j=1,...,n$ we have $P^\ttt_{j,1}(\uu)\in R^\times$. By \eqref{Eqninf1psP} and the fact that $t_j\in \mfrak_{\Acal,e}\subseteq \Ocal_{\Acal,e}\subseteq \Ocal_{A,e}$ (cf. Section \ref{SecLogExp1}), this implies that the image of $\tilde{\gamma}_m^\#$ contains an element of the form $uz+z^2h\bmod z^{m+1}$ for some $u\in R^\times$ and some $h\in K[z]$. As such an element generates $E_m$ as a $K$-algebra, we obtain that $\tilde{\gamma}_m^\#:\Ocal_{A,e}\to E_m$ is surjective. Thus, $\tilde{\gamma}_m$ is a closed immersion. 

Since $u\in R^\times$, an analogous argument works for $\tilde{\gamma}^R_m$ and $\tilde{\gamma}'_m$ if $m<p$.
\end{proof}
%%
%%


%%%%
%%%%
%%%%
\subsection{$\omega$-integrality} 
%%
%%
\begin{lemma}
\source{chabauty-coleman-bound-surfaces}
\notready\label{LemmaDiffExpl} We have $H^0(V_m, \Omega^1_{V_m/K})=(K[z]/(z^m))dz$. Furthermore, if $m\le p-2$ then $H^0(\Vcal_m,\Omega^1_{\Vcal_m/R}) = (R[z]/(z^m))dz$ and $H^0(V'_m, \Omega^1_{V'_m/\kk})=(\kk[z]/(z^m))dz$. In particular, for $m\le p-2$ we have that $H^0(\Vcal_m,\Omega^1_{\Vcal_m/R})=\Omega^1_{\Ecal_m/R} $ is a free $R$-module of rank $m$.
\source{chabauty-coleman-bound-surfaces}
\end{lemma}
\begin{proof}  Note that $H^0(V_m,\Omega^1_{V_m/K})=\Omega^1_{E_m/K}=K[z]dz/(z^{m+1}, d(z^{m+1}))$. The first claim follows from the fact that $d(z^{m+1})=(m+1)z^mdz$ and $m+1\in K^\times$. The rest is proved similarly. 
\end{proof}
%%
%%

Recall that for the additive Lie group $\pfrak$ the space of invariant differentials is $\Omega(\pfrak)=K\cdot dz$. For each $m$ we define the $K$-linear map 
$$
\rho_m:\Omega(\pfrak)\to \Omega^1_{E_m/K}=H^0(V_m, \Omega^1_{V_m/K})
$$ 
by the rule $\rho_m(dz)=dz\in \Omega^1_{E_m/K}$.

%%
%%
\begin{lemma}[Analytic-algebraic compatibility]
\source{chabauty-coleman-bound-surfaces}
\notready \label{LemmaCompatibility} Let $\gamma=(\ttt,\uu)$ be a $1$-parameter subgroup of $\Acal$ and let $m\ge 0$. The map on invariant differentials $\tilde{\gamma}^\bullet : H^0(A,\Omega^1_{A/K})\to \Omega(\pfrak)$ induced by the Lie group morphism $\tilde{\gamma}$ and the map $\tilde{\gamma}_m^\bullet : H^0(A,\Omega^1_{A/K})\to H^0(V_m,\Omega^1_{V_m/K})$ induced by the scheme morphism $\tilde{\gamma}_m:V_m\to A$ satisfy $\rho_m\circ \tilde{\gamma}^\bullet=\tilde{\gamma}_m^\bullet$.
\source{chabauty-coleman-bound-surfaces}
\end{lemma}
\begin{proof} For a $K$-algebra $B$, the module of universally finite differentials over $K$ (cf. \cite{KunzDiff} Sec. 11) is denoted by $\widetilde{\Omega}_{B/K}$. We recall that in the special case $B=K[[x_1,...,x_n]]$ we have $\widetilde{\Omega}_{B/K}=\bigoplus_{j=1}^n B\cdot dx_j$ and the map $d:B\to \widetilde{\Omega}_{B/K}$ is continuous for the $(x_1,...,x_n)$-topology.

Let $\kappa:\widetilde{\Omega}_{\widehat{\Ocal}_{A,e}/K}\to \widetilde{\Omega}_{K[[z]]/K}$ be the map induced by the continuous ring map $\widehat{\gamma}^\# : \widehat{\Ocal}_{A,e}\to K[[z]]$ given by the power series expansion of $\tilde{\gamma}$, cf. \eqref{Eqnhatgamma}. Since the map $\tilde{\gamma}^\bullet$ might be computed from the power series expansion of the Lie group morphism $\tilde{\gamma}: \pfrak\to U_e$ by taking differentials term-by-term, we deduce that the following diagram commutes:
$$
\begin{tikzcd}
  H^0(A,\Omega^1_{A/K}) \arrow{r}{\tilde{\gamma}^\bullet}   \arrow[hook]{d}   & \Omega(\pfrak) \arrow[hook]{d} \\
  \widetilde{\Omega}_{\widehat{\Ocal}_{A,e}/K} \arrow{r}{\kappa} &\widetilde{\Omega}_{K[[z]]/K}.
\end{tikzcd}
$$
The map $H^0(A,\Omega^1_{A/K})\to \widetilde{\Omega}_{\widehat{\Ocal}_{A,e}/K}$ factors through $\Omega^1_{\Ocal_{A,e}/K}$. The quotient $K[[z]]\to E_m$ induces a map $\widetilde{\Omega}_{K[[z]]/K}\to \Omega^1_{E_m/K}$ and one checks (evaluating on $dz$) that this map composed with $\Omega(\pfrak)\to \widetilde{\Omega}_{K[[z]]/K}$ is $\rho_m$. From these two observations we obtain the commutative diagram
$$
\begin{tikzcd}
   & H^0(A,\Omega^1_{A/K}) \arrow{r}{\tilde{\gamma}^\bullet}   \arrow[hook]{d}  \arrow[hook]{ld}  &  \Omega(\pfrak) \arrow[hook]{d} \arrow{rd}{\rho_m} & \\
\Omega^1_{\Ocal_{A,e}/K} \arrow[hook]{r}& \widetilde{\Omega}_{\widehat{\Ocal}_{A,e}/K} \arrow{r}{\kappa} &\widetilde{\Omega}_{K[[z]]/K}\arrow{r} & \Omega^1_{E_m/K}.
\end{tikzcd}
$$
The composition $\Omega^1_{\Ocal_{A,e}/K}\to \Omega^1_{E_m/K}$ of the bottom maps is the morphism induced by $\tilde{\gamma}^\#_m$, by \eqref{Eqngammas}. Finally, this map composed with the inclusion $H^0(A,\Omega^1_{A/K})\to \Omega^1_{\Ocal_{A,e}/K}$ is $\tilde{\gamma}^\bullet_m$, by Lemma \ref{LemmaCommLoc}.
\end{proof}
%%
%%


%%
%%
\begin{lemma}
\source{chabauty-coleman-bound-surfaces}
\notready\label{Lemmawintgamma} Let $\gamma=(\ttt,\uu)$ be a $1$-parameter subgroup of $\Acal$, let $m\ge 0$, and let $\omega \in \Hcal(\gamma)$. Then $\tilde{\gamma}_m:V_m\to A$ is $\omega$-integral. Furthermore, if $m\le p-2$ and $\omega$ extends to a section $\widetilde{\omega}\in H^0(\Acal,\Omega^1_{\Acal/R})$, then $\tilde{\gamma}'_m:V'_m\to A'$ is $\omega'$-integral, where $\omega'$ is the restriction of $\widetilde{\omega}$ to $A'$. 
\source{chabauty-coleman-bound-surfaces}
\end{lemma}
%%
%%
\begin{proof} By Lemma \ref{LemmaHisKernel} we have $\tilde{\gamma}^\bullet (\omega)=0$. Lemma \ref{LemmaCompatibility} gives $\tilde{\gamma}^\bullet_m(\omega)=\rho_m(\tilde{\gamma}^\bullet (\omega))=0$. Therefore, $\tilde{\gamma}_m:V_m\to A$ is $\omega$-integral.

Suppose now that $m\le p-2$. The morphism $\tilde{\gamma}^R_m:\Vcal_m\to \Acal$ induces a map on differentials $(\tilde{\gamma}^R_m)^\bullet : H^0(\Acal,\Omega^1_{\Acal/R})\to H^0(\Vcal_m,\Omega^1_{\Vcal_m/R})$. By base change we get a commutative diagram
$$
\begin{tikzcd}
H^0(A',\Omega^1_{A/\kk})  \arrow{d}{(\tilde{\gamma}'_m)^\bullet} &    H^0(\Acal,\Omega^1_{\Acal/R}) \arrow{l} \arrow{r} \arrow{d}{(\tilde{\gamma}^R_m)^\bullet} &   H^0(A,\Omega^1_{A/K}) \arrow{d}{\tilde{\gamma}_m^\bullet} \\
H^0(V'_m,\Omega^1_{V'_m/\kk})   &   H^0(\Vcal_m,\Omega^1_{\Vcal_m/R}) \arrow{l} \arrow{r}& H^0(V_m,\Omega^1_{V_m/K})
\end{tikzcd}
$$

Since $H^0(\Vcal_m,\Omega^1_{\Vcal_m/R})$ is free as an $R$-module (cf. Lemma \ref{LemmaDiffExpl}), the base change morphism $H^0(\Vcal_m,\Omega^1_{\Vcal_m/R}) \to H^0(V_m,\Omega^1_{V_m/K})$  is injective. Hence, $(\tilde{\gamma}^R_m)^\bullet(\widetilde{\omega})=0$ because $\tilde{\gamma}^\bullet_m(\omega)=0$ and the right square of the previous diagram commutes.

By commutativity of the left square of the previous diagram, we deduce that $(\tilde{\gamma}'_m)^\bullet(\omega')=0$. Hence $\tilde{\gamma}'_m:V'_m\to A'$ is $\omega'$-integral.
\end{proof}
%%
%%





%%%%%%%%%%%%%%%%%%%%%%%%%%%%%%%%%%%%%%
%%%%%%%%%%%%%%%%%%%%%%%%%%%%%%%%%%%%%%
%%%%%%%%%%%%%%%%%%%%%%%%%%%%%%%%%%%%%%
%%%%%%%%%%%%%%%%%%%%%%%%%%%%%%%%%%%%%%
%%%%%%%%%%%%%%%%%%%%%%%%%%%%%%%%%%%%%%
%%%%%%%%%%%%%%%%%%%%%%%%%%%%%%%%%%%%%%
%%%%%%%%%%%%%%%%%%%%%%%%%%%%%%%%%%%%%%
%%%%%%%%%%%%%%%%%%%%%%%%%%%%%%%%%%%%%%
%%%%%%%%%%%%%%%%%%%%%%%%%%%%%%%%%%%%%%
%%%%%%%%%%%%%%%%%%%%%%%%%%%%%%%%%%%%%%
%%%%%%%%%%%%%%%%%%%%%%%%%%%%%%%%%%%%%%
%%%%%%%%%%%%%%%%%%%%%%%%%%%%%%%%%%%%%%
%%%%%%%%%%%%%%%%%%%%%%%%%%%%%%%%%%%%%%
%%%%%%%%%%%%%%%%%%%%%%%%%%%%%%%%%%%%%%
%%%%%%%%%%%%%%%%%%%%%%%%%%%%%%%%%%%%%%
%%%%%%%%%%%%%%%%%%%%%%%%%%%%%%%%%%%%%%
%%%%%%%%%%%%%%%%%%%%%%%%%%%%%%%%%%%%%%
%%%%%%%%%%%%%%%%%%%%%%%%%%%%%%%%%%%%%%



%%%%%%%%%%%%%%%%%%%%%%%%%%%%%%%%%%%%%%
%%%%%%%%%%%%%%%%%%%%%%%%%%%%%%%%%%%%%%
%%%%%%%%%%%%%%%%%%%%%%%%%%%%%%%%%%%%%%
%%%%%%%%%%%%%%%%%%%%%%%%%%%%%%%%%%%%%%
%%%%%%%%%%%%%%%%%%%%%%%%%%%%%%%%%%%%%%
%%%%%%%%%%%%%%%%%%%%%%%%%%%%%%%%%%%%%%
